\section{From the linear reserve to the sharp rate}
\label{upgrade:section}

The local reserve inequality does not by itself give the sharp rate
$F(I)$ at every state. We now join their maximum
\[
 \widehat{\mathcal B}_m(I)
 =\max\{\mathcal B_m(I),F(I)\}.
\]
Closure under a maximum is a useful feature of the local-production
framework \cite[Lemma~3]{ChenGohariNair2025}, but $F$ alone does not
satisfy a joining inequality for every law. We therefore prove the
needed comparison in the part of the domain where the reserve is
insufficient. Two comparisons organize the proof. Above half information, the
reserve already dominates $F$. Below half information, the demand
from $F$ increases with the common mean of the child states. Thus
every reflected core can be paid at its one-corner wall by the same
thinning theorem used for a one-corner optimizer.

\subsection{The reserve above half information}

\begin{lemma}[The high-information reserve]
\label{upgrade:high-information}
For $|m|<1$ and $L/2\le I<H(m)$,
\begin{equation}
 \mathcal B_m(I)\ge F(I).
 \label{upgrade:high-information-bound}
\end{equation}
The inequality is strict if $m\ne0$. Consequently a state with
$F(I)>\mathcal B_m(I)$ must have $I<L/2$.
\end{lemma}

\begin{proof}
Recall $S(I)=F(I)-2I$ on $[0,L)$ and put
\[
 k_0=L-\frac12,\qquad
 \Xi(I)=(I-k_0)S'(I)-S(I).
\]
Since $S(0)=S'(0)=0$ and $S''=F''$ on this interval,
\begin{equation}
 \Xi(I)=\int_0^I(s-k_0)F''(s)\,ds.
 \label{upgrade:centered-curvature}
\end{equation}
Lemma~\ref{scalar:curvature-growth} says that
$w(s)=e^{-2s}F''(s)$ strictly increases. Also
$0<k_0<L/2$, and the binary identity $e^L=2$ gives the exact
balance
\begin{equation}
 \int_0^{L/2}(s-k_0)e^{2s}\,ds
 =\left[\frac12e^{2s}(s-L)\right]_0^{L/2}=0.
 \label{upgrade:exponential-balance}
\end{equation}
On the part of the integral below $k_0$, the factor $s-k_0$ is
negative and $w(s)\le w(k_0)$. Above $k_0$, one has
$s-k_0>0$ and $w(s)\ge w(k_0)$. Replacing $w(s)$ by $w(k_0)$
therefore gives a lower
bound on both parts, strictly so on intervals of positive length.
Equations \eqref{upgrade:centered-curvature}--
\eqref{upgrade:exponential-balance} imply
\[
 \Xi(L/2)
 =\int_0^{L/2}(s-k_0)e^{2s}w(s)\,ds>0.
\]
Since $\Xi'(I)=(I-k_0)F''(I)>0$ for $I\ge L/2$, we have
$\Xi(I)>0$ throughout that range.

For the mean-dependent comparison, set
\[
 \alpha=1-m^2,\qquad
 c_m=H(m)-\alpha L=Lm^2-\psi(m).
\]
The elementary entropy bounds give
$0\le c_m\le m^2k_0$. Hence, for $I\ge L/2$, the argument
\[
 s=L-\frac{H(m)-I}{\alpha}
   =\frac{I-c_m}{\alpha}
\]
is positive; it is less than $L$ because $I<H(m)$. No clipping
occurs in this comparison. The tangent inequality for the convex
function $S$, and $S'\ge0$, now yield
\begin{align*}
 \mathcal B_m(I)-F(I)
 &=\alpha S(s)-S(I)\\
 &\ge\alpha[S(I)+(s-I)S'(I)]-S(I)\\
 &=m^2[IS'(I)-S(I)]-c_mS'(I)\\
 &\ge m^2[(I-k_0)S'(I)-S(I)]
 =m^2\Xi(I)\ge0.
\end{align*}
The inequality is strict at a nonzero mean. At mean zero,
$\alpha=1$, $c_m=0$, and the two profiles agree identically.
\end{proof}

\subsection{A direct comparison with the one-corner wall}

\begin{lemma}[The demand below half information]
\label{wall:monotonicity}
Let $m,d>0$, $m+d<1$, and
\[
 0<e\le\bar h(m,d):=\frac{H(m+d)+H(m-d)}2.
\]
Put $I=H(m)-e$, $j=\bar h(m,d)-e$, and $J=I-j$. If
$I\le L/2$, then at fixed $d,e$,
\begin{equation}
 \partial_m\{F(I)-F(j)\}>0.
 \label{wall:positive-derivative}
\end{equation}
At $j=0$, the derivative is interpreted from the physical side.
\end{lemma}

\begin{proof}
We compare the rate slopes with the entropy slopes. First,
\begin{equation}
 \frac{F''(u)}{F'(u)}<2\qquad(0\le u\le L/2).
 \label{wall:rate-slope}
\end{equation}
Indeed, set $h=L-u\ge L/2$. The lower bound
$F''(u)\ge1/[2(L-u)^2]$ from
Lemma~\ref{scalar:scaled-curvature} gives
\[
 F'(u)\ge2+\frac12\left(\frac1h-\frac1L\right).
\]
Using the upper bound from the same lemma, we get
\[
 2h^2F'(u)\ge(4-L^{-1})h^2+h
 \ge L^2+\frac L4>\frac43L^2\ge h^2F''(u).
\]
The quadratic in $h$ increases for $h>0$, and the strict
comparison uses $L<3/4$. This proves \eqref{wall:rate-slope}.
Since $0\le j<I\le L/2$, integration gives
\begin{equation}
 \log\frac{F'(I)}{F'(j)}<2(I-j)=2J.
 \label{wall:rate-ratio}
\end{equation}

For the entropy slopes, write
\[
 A=-\partial_m\bar h(m,d)
   =\frac{V(m+d)+V(m-d)}2>0.
\]
We claim the reverse comparison
\begin{equation}
 \log\frac{A}{V(m)}\ge2J.
 \label{wall:entropy-ratio}
\end{equation}
Hold $m>0$ fixed and set $v_\pm=V(m\pm d)$. The function
$G(v)=\cosh^2v-v^2$ is even and strictly increasing for $v>0$,
because $G'(v)=\sinh(2v)-2v>0$. Direct differentiation gives
\[
 \partial_d\left\{\log\frac{A}{V(m)}-2J\right\}
 =\frac{G(v_+)-G(v_-)}{v_++v_-}>0.
\]
Indeed, $A_d=(\cosh^2v_+-\cosh^2v_-)/2$ and
$2J_d=v_+-v_-$. The numerator and denominator are positive
since $m+d>|m-d|$, so $v_+>|v_-|$. At $d=0$ the expression
in braces is zero. Integration proves \eqref{wall:entropy-ratio}.

Finally $I_m=-V(m)$ and $j_m=-A$. Therefore
\[
 \partial_m\{F(I)-F(j)\}
 =-F'(I)V(m)+F'(j)A>0
\]
by \eqref{wall:rate-ratio}--\eqref{wall:entropy-ratio}.
The regular right derivative $F'(0)=2$ includes a capacity endpoint.
\end{proof}

\subsection{The local inequality and dimension induction}

\begin{theorem}[Joining the strengthened profile]
\label{upgrade:joining}
Let $(A,B)$ have a joint law on $[-1,1]^2$ with
$e_a=\E H(A)>0$ and $e_b=\E H(B)>0$. Put
\[
 a=\E A,\quad b=\E B,\quad
 m=\frac{a+b}{2},\quad e=\frac{e_a+e_b}{2},\quad
 I=H(m)-e,
\]
and $I_a=H(a)-e_a$, $I_b=H(b)-e_b$. Then
\begin{equation}
 \E c(A,B)\ge
 \widehat{\mathcal B}_m(I)
 -\frac{\widehat{\mathcal B}_a(I_a)
             +\widehat{\mathcal B}_b(I_b)}2.
 \label{upgrade:local-inequality}
\end{equation}
\end{theorem}

\begin{proof}
An infinite joining cost makes the assertion immediate. Otherwise
positive child entropies give $|a|,|b|<1$, and entropy concavity
makes all the profile arguments physical.

If the parent maximum is attained by $\mathcal B_m(I)$, apply
the local reserve inequality of Theorem~\ref{reserve:joining}
and enlarge the two child terms to their
$\widehat{\mathcal B}$ values. This includes every tie at the
parent.

It remains to suppose $F(I)>\mathcal B_m(I)$. Put
$d=|a-b|/2$ and
\[
 \bar h(m,d)=\frac{H(m+d)+H(m-d)}2,
 \qquad j=\bar h(m,d)-e=\frac{I_a+I_b}{2}.
\]
Since $\widehat{\mathcal B}\ge F$, convexity at the original
child states gives
\begin{equation}
 \frac{\widehat{\mathcal B}_a(I_a)
             +\widehat{\mathcal B}_b(I_b)}2
 \ge\frac{F(I_a)+F(I_b)}2\ge F(j).
 \label{upgrade:original-child-pooling}
\end{equation}
For $d=0$, one has $I=j$ and nonnegativity of cost proves the
claim. Otherwise reflect to $m>0$: mean zero is impossible in
this strict-parent case, because $\mathcal B_0=F$.
Let
\[
 J_0(x,d,e)=F(H(x)-e)-F(\bar h(x,d)-e).
\]
The actual law is feasible in the common-budget cost problem, so
it is enough to prove
\begin{equation}
 \mathcal K(m,d,e)\ge J_0(m,d,e).
 \label{upgrade:pure-increment}
\end{equation}

Use the exact optimizer from
Proposition~\ref{canonical:coverage}. In its one-positive-corner
geometry, the entropy is at least the transition value
\eqref{canonical:transition}; Theorem~\ref{thinning:high-value}
therefore proves \eqref{upgrade:pure-increment}.

In the reflected geometry \eqref{canonical:reflected}, write
\[
 d=kr,\qquad e=kH(r),\qquad \ell=1-k,
 \qquad 0<r<1,\quad0<k\le1.
\]
Then $0<m\le\ell$, and the cost is $dV(r)$. Hold $d,e,k,r$
fixed and increase the common mean from $m$ to $\ell$.
The law with mass $k$ at $(r,-r)$ and masses
$(1-k+x)/2$, $(1-k-x)/2$ at the equal corners makes every
$x\in[m,\ell]$ physical at the same entropy and cost.
Its means remain interior, since
$\ell+d=1-k(1-r)<1$.

Lemma~\ref{upgrade:high-information} gives $H(m)-e<L/2$.
The function $H(x)-e$ decreases along this interval, so
Lemma~\ref{wall:monotonicity} applies throughout. Consequently
\[
 J_0(m,d,e)\le J_0(\ell,d,e).
\]
At $x=\ell$, the negative-corner mass is zero, and the state is
at the one-corner transition. Theorem~\ref{thinning:high-value}
gives
\[
 J_0(\ell,d,e)\le dV(r)=\mathcal K(m,d,e).
\]
This proves \eqref{upgrade:pure-increment}. A zero-length
interval requires only the last comparison. The case $k=1$
would force $m=0$ and is already covered by the reserve branch.

Only the information bound is retained while the mean increases;
no claim is made about which profile stays active there. The child
terms in \eqref{upgrade:original-child-pooling} always use the
original entropies $e_a,e_b$, not the entropies of a replacement
optimizer. Together with that pooling inequality, the cost bound
completes \eqref{upgrade:local-inequality}.
\end{proof}

\begin{proof}[Proof of Theorem~\ref{upgrade:production}]
Induct on dimension. At dimension zero, a constant field has
$D=I=0$ and $\widehat{\mathcal B}_m(0)=0$. For a coordinate
split $g_+,g_-$, the exact energy decomposition is
\[
 D(g)=\frac{D(g_+)+D(g_-)}2+\E c(g_+,g_-).
\]
Each within-face term has its face probability $1/2$; symmetry
of the joining cost combines the two orientations of every
joining edge into the last expectation.
The induction hypothesis bounds the within-face productions,
and Theorem~\ref{upgrade:joining} supplies the remaining
difference on the joining edges. Their sum is the parent
$\widehat{\mathcal B}$ value. Every face is interior-valued,
so all child entropies are positive and all local uses are physical.
\end{proof}
